\chapter{Le basi dello strong gravitational lensing}
\label{cap:SGL}

\section{Geometria e parametri fondamentali}
\subsection{Il principio di Fermat}
\subsection{L'equazione della lente e l'angolo di deflessione}

\section{Il potenziale della lente}
\subsection{Convergenza}
\subsection{Ingrandimento e distorsione}
\subsection{Formazione delle immagini}
\subsection{Il modello isotermo ellittico}

\section{Le galassie ellittiche come lenti}
\subsection{Vantaggi delle galassie ellittiche massicce come lenti forti}
%  Sez. 1.1 e 1.2 di Shajib, Strong Lensing by Galaxies
\subsection{Il profilo di massa totale}
% Sez. 4.1 di Shajib, Strong Lensing by Galaxies
\subsection{Osservabili e degenerazioni principali}
% Sez. 3.3.2 di Shajib, Strong Lensing by Galaxies  (degenerazioni nei modelli parametrici, definizione della MSD, necessità di combinare lensing e cinematica per romperla)